
\subsection*{1.3.}
Rappelons (IV 2.10) que le bifoncteur
@@M0@@
où @@M1@@ est un @@M2@@-module localement libre et @@M3@@ un @@M4@@-module cohérent, induit une loi de composition
@@M5@@
qui fait de @@M6@@ un @@M7@@-module unitaire. Le résultat-clé du présent exposé est le

\textbf{Théorème 1.3.2.} \uline{Soit} @@M8@@ \uline{un schéma noethérien. Pour tout} @@M9@@ \uline{la loi de composition} (1.3.1) \uline{induit par restriction une loi de composition}
@@M10@@

\textbf{Corollaire 1.3.3.} \uline{Si} @@M11@@ \uline{est un schéma noethérien de dimension} @@M12@@\uline{, alors l'homomorphisme canonique}
@@M13@@
\uline{défini par} @@M14@@ (@@M15@@\uline{) induit pour tout} @@M16@@ \uline{un homomorphisme}
@@M17@@

\emph{Démonstration :} Montrons d'abord que (1.3.3), pour @@M18@@ et pour \emph{tout} sous-schéma intègre fermé de dimension @@M19@@, implique (1.3.2) pour @@M20@@. Grâce à (1.1.4) il suffit de montrer que pour tout @@M21@@ et pour \emph{tout} sous-schéma intègre fermé @@M22@@ de dimension @@M23@@ on a
